\chapter{Paroxysmal Nocturnal Hemoglobinuria Test}

\section*{What Is This Test?}
Paroxysmal nocturnal hemoglobinuria (PNH) testing detects blood cells that lack proteins anchored to the cell membrane by glycosylphosphatidylinositol (GPI). The abnormal clone results from an acquired mutation in \textit{PIGA} and can cause intravascular hemolysis, bone-marrow failure, and thrombosis.

\section*{Alternative Names}
\begin{itemize}
  \item PNH Test
  \item PNH Flow Cytometry
  \item GPI-Deficient Cell Testing
  \item FLAER Assay
\end{itemize}
\section*{How It Is Performed}
Flow cytometry analyzes red cells and white cells using fluorescent aerolysin (FLAER) and antibodies to GPI-linked proteins such as CD55 and CD59. Testing multiple cell lineages improves sensitivity because red-cell populations may be reduced by hemolysis or recent transfusion.

\section*{Why It Is Ordered}
\begin{itemize}
  \item Investigation of unexplained intravascular hemolysis, hemoglobinuria, or elevated LDH
  \item Evaluation of aplastic anemia, myelodysplastic disease, or unexplained cytopenias
  \item Assessment of unusual or recurrent thrombosis, especially at abdominal or cerebral venous sites
  \item Monitoring a known PNH clone during specialist follow-up
\end{itemize}

\section*{Interpretation and Limitations}
The report usually gives the size of the PNH clone in each tested lineage. A small clone may be clinically significant in bone-marrow-failure disorders, while clone size alone does not predict symptoms or thrombosis risk. Recent transfusion, severe leukopenia, and inadequate specimen handling can reduce sensitivity. Results require correlation with blood counts, reticulocytes, LDH, bilirubin, haptoglobin, and clinical findings.
\section*{Specimen and Preparation}
Testing uses fresh EDTA blood and requires no fasting. Recent transfusion, severe leukopenia, hemolysis, or delayed transport can reduce the measurable clone, so provide transfusion dates and current treatment. Testing more than one cell lineage improves sensitivity because red cells may be replaced or lost.

\section*{Risks and Safety}
Physical risk is limited to venipuncture. Suspected PNH can be urgent when there is dark urine, severe anemia, abdominal pain, breathlessness, headache, or unusual thrombosis. A small clone does not necessarily mean mild disease, and a result should not be used without assessment of hemolysis and marrow function.

\section*{Understanding Results and Follow-up}
The report should state the percentage and assay sensitivity for each tested lineage and whether GPI-linked proteins or FLAER were assessed. Follow-up may include CBC, reticulocytes, LDH, bilirubin, haptoglobin, renal assessment, marrow evaluation, and specialist review. Clone size is interpreted with hemolysis, cytopenias, thrombosis, and symptoms rather than as a stand-alone severity score.

\section*{Regulatory Status and Authorization Date}
This chapter describes a test category, not a specific commercial product. FDA, CDSCO, EU IVDR, and MHRA approval or registration dates therefore cannot be assigned without the assay manufacturer, product name, instrument, and intended use. For laboratory-developed tests, an individual agency approval date may not exist. See the Preface for the interpretation of regulatory status.

\clearpage
